% IV. Програмна реализация. Изходният текст още не е написан: този раздел описва
% замисъла на реализацията и се допълва с листинги, когато кодът съществува.
\section{Програмна реализация}
\label{sec:implementation}

\subsection{Организация на изходния текст}

Хранилището разделя изходния текст на три части. Заглавните файлове с публичните договори
стоят в \code{include/}, реализацията в \code{src/}, а клиентската част в \code{web/}.
Тестовете са в \code{tests/}. Изграждането се управлява от CMake, което позволява
проектът да се компилира с различни компилатори без промяна в описанието на изграждането.
\TODO{дърво на хранилището с кратко описание на всеки файл, след като реализацията
приключи}

\subsection{Мрежов слой на сървъра}

Мрежовият слой обработва два протокола върху един и същ слушащ сокет. Първата заявка по
всяка връзка е обикновена HTTP/1.1 заявка~\cite{rfc9112}. Ако тя носи заглавните части за
надграждане към WebSocket, слоят изчислява отговора на ръкостискането по правилото от
спецификацията~\cite{rfc6455} и предава връзката на кадровия анализатор. В противен случай
заявката се обслужва като заявка за статичен ресурс и връзката се затваря или се запазва
според семантиката на HTTP~\cite{rfc9110}.

Кадровият анализатор на WebSocket е мястото, където проектът е най-чувствителен към
грешки, защото работи с данни, дошли от мрежата. Дължината на полезния товар се чете от
кадъра, тоест от ненадежден източник, и затова се проверява срещу горна граница, преди да
бъде заделена памет. Кадрите от клиент към сървър задължително носят маска, а кадър без
маска се отхвърля вместо да се обработва. \TODO{листинг на анализатора на кадри с коментар
на изходния текст}

\subsection{Сериализатори}

Двата сериализатора реализират един и същ вътрешен договор: приемат пакет със стойности и
връщат низ за изпращане. Реализацията за нотацията за обекти~\cite{rfc8259} се грижи за
екранирането на специалните знаци в низовете. Реализацията за разширяемия език за
разметка~\cite{xml10} се грижи за екранирането на знаците със специално значение в
разметката и за коректното изписване на празни елементи. И двете се проверяват със
самостоятелни тестове, защото грешка в екранирането не се вижда в обичайния случай и се
проявява само при определени входни данни.

\subsection{Клиентска част}

Клиентската част е един документ, един стилов файл и няколко скриптови модула без външни
зависимости. Модулът за връзката отговаря за установяването на канала и за повторното
свързване след прекъсване. Модулът за състоянието прилага пристигналите пакети върху
текущата картина. Модулът за свързване на данни изпълнява описания в
Раздел~\ref{sec:design} списък от двойки. Модулът за диаграмите изгражда и обновява
структурната графика~\cite{svg2}.

Обработката на събития е съсредоточена на едно място чрез делегиране: обработчикът се
закача на общ родителски възел и разпознава източника на събитието по неговите данни, а не
по отделен обработчик на всеки възел. Това решение произтича пряко от факта, че възлите на
диаграмата се създават и унищожават при всяко обновяване. \TODO{листинги на клиентските
модули и екранни снимки на готовия интерфейс}
